% Slide "Several cores: speedup over plain scalar", left panel: f32. Speedup over the plain scalar
% build of the same compiler at the same core count, lid-driven cavity (SweepgenLdcCPU), one 64^3
% block per MPI process, best of 3 runs of 200 steps. One session (board, scripts/presentation4/all.sh,
% logs docs/presentation4/board-logs/board/ldc-{gcc,clang}-{scalar,autovec,rvv}-f32-{1,4,8}core.log),
% printed by
%   python3 scripts/presentation4/slide_numbers.py --tex docs/presentation4/board-logs/board
% Bars and numbers: GCC 14.2. Diamonds: Clang 18.1, against its own plain scalar (plain scalar = 1
% for both compilers, so no diamond there). Series as figures/mlups-multicore.tex: plain scalar =
% scalar, best scalar = scalar_cse_aa, auto-vec. in place = autovec_cse_aa (hatched),
% RVV in place + prefetch = explicit_cse_aa_pf, RVV split loop = explicit_cse_aa_split (in place,
% no software prefetch: 0 hints).
% Styles mc, pl, bs, av, rv, sp, cl: preamble of slides.tex. Same ymax as speedup-multicore-f64.tex.
% presentation5: labels with two digits (notes5 F8).
\begin{tikzpicture}
  \begin{axis}[mc, ymax=7.5, title={f32}, ylabel={speedup over plain scalar}]
    \addplot[pl] coordinates {(1,1)[1.0] (4,1)[1.0] (8,1)[1.0]};
    \addplot[bs] coordinates {(1,1.447)[1.4] (4,1.384)[1.4] (8,1.425)[1.4]};
    \addplot[av] coordinates {(1,1.449)[1.4] (4,1.379)[1.4] (8,1.427)[1.4]};
    \addplot[rv] coordinates {(1,4.738)[4.7] (4,3.397)[3.4] (8,2.211)[2.2]};
    \addplot[sp] coordinates {(1,5.548)[5.5] (4,3.472)[3.5] (8,2.222)[2.2]};
    \addplot[cl,xshift=-6pt] coordinates {(1,1.554) (4,1.466) (8,1.513)};
    \addplot[cl,xshift=0pt] coordinates {(1,1.542) (4,1.466) (8,1.511)};
    \addplot[cl,xshift=6pt] coordinates {(1,4.551) (4,3.290) (8,2.200)};
    \addplot[cl,xshift=12pt] coordinates {(1,5.550) (4,3.425) (8,2.250)};
  \end{axis}
\end{tikzpicture}
